% ===================================================================
\section{The Kernel Category}
\label{ch:kernel}
% ===================================================================

\subsection{Overview}

FRFP is built on a minimal three-object category $\Czero$ whose structure encodes
the fundamental distinction between explicit (machine-manipulable) and tacit
(human correctness-bearing) computation. This chapter defines $\Czero$, its six
primitive morphisms, their typing rules, and the core theorems that follow directly
from the definitions --- without requiring any additional axioms.

% -------------------------------------------------------------------
\subsection{Objects}
% -------------------------------------------------------------------

\begin{definition}[Kernel category objects]
The kernel category $\Czero$ has exactly three objects:
\[
  \mathrm{Ob}(\Czero) = \{\varnothing,\ \Eobj,\ \Tobj\}
\]
\end{definition}

\begin{center}
\begin{tabular}{lll}
\toprule
Object & Name & Interpretation \\
\midrule
$\varnothing$ & Empty & The absent or uninitialized explicit state \\
$\Eobj$ & Explicit object & Machine-manipulable representations \\
$\Tobj$ & Tacit object & Human correctness-bearing knowledge \\
\bottomrule
\end{tabular}
\end{center}

The object set is closed. No extensions to $\Czero$ are permitted.
\leanref{Frfp.Core.Kernel.Object}

The objects are sometimes written with subscript notation: $\Eobj$ (the explicit
\emph{object}) is distinct from $\Ecat$ (the \emph{category} of explicit
morphisms). This distinction is enforced in the Lean formalization by separate
types.

% -------------------------------------------------------------------
\subsection{Primitive Morphisms}
% -------------------------------------------------------------------

\begin{definition}[Primitives]
The six primitive morphisms of $\Czero$ are:
\end{definition}

\begin{center}
\begin{tabular}{llll>{\raggedright\arraybackslash}p{4.5cm}}
\toprule
Primitive & Source & Target & Name & Role \\
\midrule
$\RI$ & $\varnothing$ & $\Eobj$ & Representation Initiation & Creates an explicit representation \\
$\EC$ & $\Eobj$ & $\Eobj$ & Explicit Computation & AI processing within explicit space \\
$\ED$ & $\Eobj$ & $\Eobj$ & Explicit Diagnostics & AI checking/verification within explicit space \\
$\RB$ & $\Eobj$ & $\Tobj$ & Representational Backflow & The unique boundary crossing morphism \\
$\TE$ & $\Tobj$ & $\Tobj$ & Tacit Evaluation & Human evaluation within tacit space \\
$\HFD$ & $\Tobj$ & $\Tobj$ & Human Final Decision & Human closure within tacit space \\
\bottomrule
\end{tabular}
\end{center}

\leanref{Frfp.Core.Kernel.Primitive}

The primitive set is a closed enumeration. No new primitives may be added without
a complete re-audit of all downstream theorems.

% -------------------------------------------------------------------
\subsection{Subcategories}
% -------------------------------------------------------------------

\begin{definition}[Explicit subcategory]
$\Ecat$ is the full subcategory of morphisms targeting $\Eobj$:
\[
  \Ecat = \{p : \mathrm{Primitive} \mid p.\mathrm{target} = \Eobj\}
        = \{\RI, \EC, \ED\}
\]
\end{definition}

\begin{definition}[Tacit subcategory]
$\Tcat$ is the full subcategory of morphisms targeting $\Tobj$:
\[
  \Tcat = \{p : \mathrm{Primitive} \mid p.\mathrm{target} = \Tobj\}
        = \{\RB, \TE, \HFD\}
\]
\end{definition}

\leanref{Frfp.Core.Kernel.Ecat}, \leanref{Frfp.Core.Kernel.Tcat}

The partition $\Ecat \cup \Tcat$ covers all six primitives and is disjoint.

% -------------------------------------------------------------------
\subsection{Core Typing Theorems}
% -------------------------------------------------------------------

The following theorems are proved directly from the primitive definitions --- no
additional axioms are required.

\begin{theorem}[Explicit--Tacit Separation]
There is no primitive morphism from $\Tobj$ to $\Eobj$:
\[
  \forall p : \mathrm{Primitive},\
  \neg(p.\mathrm{source} = \Tobj \wedge p.\mathrm{target} = \Eobj)
\]
\end{theorem}
\begin{proof}
By case analysis on the six primitives. None has source $\Tobj$ and target
$\Eobj$. \qed
\end{proof}
\leanref{Frfp.Core.Kernel.no\_morphism\_T0\_to\_E0}

\begin{remark}
In the original published theory, ETS was listed as a primitive axiom. The Lean
formalization shows it is a consequence of the primitive typing. See
Chapter~\ref{ch:basis}.
\end{remark}

\begin{theorem}[Uniqueness of the boundary morphism]
$\RB$ is the unique primitive crossing from $\Eobj$ to $\Tobj$:
\[
  \forall p : \mathrm{Primitive},\
  (p.\mathrm{source} = \Eobj \wedge p.\mathrm{target} = \Tobj) \Rightarrow p = \RB
\]
\end{theorem}
\begin{proof}
By case analysis on the six primitives. Only $\RB$ has source $\Eobj$ and target
$\Tobj$. \qed
\end{proof}
\leanref{Frfp.Core.Kernel.RB\_unique\_boundary}

\begin{theorem}[AI executability partition]
Define \texttt{is\_AI\_executable} as $\{\RI, \EC, \ED\}$. Then:
\[
  \forall p,\
  \textit{is\_AI\_executable}(p)
  \Rightarrow (p.\mathrm{source} \in \{\varnothing, \Eobj\})
              \wedge p.\mathrm{target} = \Eobj
\]
\end{theorem}
\leanref{Frfp.Core.Kernel.AI\_preserves\_explicit}

\begin{corollary}[Human-exclusive tacit access]
$\RB$, $\TE$, and $\HFD$ are never AI-executable and always operate in or
enter tacit space.
\leanref{Frfp.Core.Kernel.is\_human\_only}
\end{corollary}

% -------------------------------------------------------------------
\subsection{Pipelines}
% -------------------------------------------------------------------

\begin{definition}[Pipeline]
A pipeline $\Pi$ is a morphism in $\Czero$: a triple $(\mathrm{src},
\mathrm{tgt}, p)$ where $p$ is the underlying primitive.
\end{definition}

\begin{center}
\begin{tabular}{llll}
\toprule
Pipeline type & Source & Target & Underlying primitives \\
\midrule
Explicit pipeline & $\varnothing$ or $\Eobj$ & $\Eobj$ & $\RI$, $\EC$, $\ED$ \\
Boundary pipeline & $\Eobj$ & $\Tobj$ & $\RB$ (unique) \\
Tacit pipeline & $\Tobj$ & $\Tobj$ & $\TE$, $\HFD$ \\
\bottomrule
\end{tabular}
\end{center}

\begin{definition}[Pipeline composition]
For pipelines $\Pi_1 : X \to Y$ and $\Pi_2 : Y \to Z$, their composition
$\Pi_2 \circ \Pi_1 : X \to Z$ is defined when
$\Pi_1.\mathrm{target} = \Pi_2.\mathrm{source}$.
\leanref{Frfp.Core.Kernel.pipeline\_compose}
\end{definition}

\begin{theorem}[Composition preserves typing]
If $\Pi_1 : X \to Y$ and $\Pi_2 : Y \to Z$, then
$(\Pi_2 \circ \Pi_1).\mathrm{source} = X$ and
$(\Pi_2 \circ \Pi_1).\mathrm{target} = Z$.
\leanref{Frfp.Core.Kernel.pipeline\_compose}
\end{theorem}

% -------------------------------------------------------------------
\subsection{Primitive Phase Partition}
% -------------------------------------------------------------------

\begin{definition}[Primitive phase]
The phase of a primitive is defined by:
\[
  \mathrm{phase}(p) =
  \begin{cases}
    \text{explicit}  & \text{if } p \in \{\RI, \EC, \ED\} \\
    \text{boundary}  & \text{if } p = \RB \\
    \text{tacit}     & \text{if } p \in \{\TE, \HFD\}
  \end{cases}
\]
\leanref{Frfp.Minimal.BasisSweep.phaseOf}
\end{definition}

\begin{theorem}[Phase partition]
The phase function is total: every primitive has exactly one phase.
\leanref{Frfp.Minimal.BasisSweep.phase\_partition}
\end{theorem}

\begin{theorem}[Boundary uniqueness by phase]
A primitive has boundary phase if and only if it is $\RB$.
\leanref{Frfp.Minimal.BasisSweep.phase\_boundary\_iff}
\end{theorem}

% -------------------------------------------------------------------
\subsection{The No-Return Principle}
% -------------------------------------------------------------------

\begin{theorem}[No tacit-to-explicit return]
There is no sequence of primitives that goes from $\Tobj$ to $\Eobj$:
\[
  \nexists\ p_1, \ldots, p_n : \mathrm{Primitive},\quad
  p_1.\mathrm{source} = \Tobj \wedge p_n.\mathrm{target} = \Eobj
\]
\end{theorem}

This follows from Theorem~1.1 (ETS) by induction: every primitive either stays
in tacit space ($\Tobj \to \Tobj$) or goes to $\Eobj$ from $\varnothing$ or
$\Eobj$.

The one-directional boundary
$\Eobj \xrightarrow{\RB} \Tobj$ is a genuine asymmetry in the categorical
structure: not a functor, not an equivalence, but a one-way transition that
formalizes the tacit grounding of AI outputs by human judgment.

% -------------------------------------------------------------------
\subsection{Immutability of the Kernel}
% -------------------------------------------------------------------

\begin{note}[Kernel immutability]
The definitions of \texttt{Object}, \texttt{Primitive},
\texttt{Primitive.source}, and \texttt{Primitive.target} are immutable. No module
in the library may shadow, redefine, or extend them.
\end{note}

This is enforced architecturally: \texttt{Kernel.lean} has no imports other than
Lean's standard library, and all other modules depend on it. Any change to
Kernel invalidates all 3305 downstream build jobs.
\leanref{Frfp.Core.Kernel --- no extends, no shadows}
